\documentclass[11pt]{article}
\usepackage[a4paper,margin=28mm]{geometry}
\usepackage{amsmath,amssymb,amsthm,hyperref}
\newcommand{\OK}{\mathcal O_K}
\newcommand{\Q}{\mathbb Q}
\newcommand{\hF}{h_{\mathrm F}}
\title{Independent verification of the differential-evaluation height bound}
\author{}\date{9 November 2026}
\begin{document}\maketitle

The evaluation argument in
\nolinkurl{modular_height_evaluation_candidate.tex} was checked
independently. Its evaluation-lattice bound, finite-discriminant
estimate, and application to the modular cover are valid under the
stated hypotheses. The final contact-height conclusion additionally
uses the displayed inherited modular realization; the existence and
numerical bounds of that realization are not reproved in this audit.

\section{The integral evaluation map}
Write $A=\OK[t]$ and let $B$ be its normalization in the curve's
function field. A normal surface is Cohen--Macaulay, and the base
surface is regular. Thus the finite normalization is flat and its
affine algebra $B$ is finite projective over $A$. This is also
explicitly stated in Javanpeykar, Section~4.2.

Let $\mathcal X$ be a regular semistable model and
$E=H^0(\mathcal X,\omega_{\mathcal X/\OK})$. A common regular
model dominates $\mathcal X$ and the normalization $\mathcal Y$.
Blowing up closed points does not change the module of global
relative canonical sections. At every codimension-one point of
$\mathcal Y$, the common model is an isomorphism locally at the
generic point of the corresponding divisor. Hence its canonical
sections are sections of $\omega_{\mathcal Y/\OK}$ in codimension
one. The latter module is reflexive, since $\mathcal Y$ is a
normal Cohen--Macaulay surface, so these sections extend globally.
Invertibility of this dualizing sheaf is unnecessary.

Finite duality on the affine chart gives
\[
 \omega_{B/\OK}=\operatorname{Hom}_A(B,A)\,dt.
\]
Because $B$ is projective, reduction modulo $t-2$ commutes with
this dual. With $C=B/(t-2)B$, evaluation therefore maps
\[
 E\longrightarrow\operatorname{Hom}_{\OK}(C,\OK).
\]
On the generic split fibre the trace functional attached to a tuple
$(f_1,\ldots,f_d)$ is $(b_i)\mapsto\sum_i f_ib_i$. Consequently
the tuple attached to $\omega$ is precisely
$(\omega/dt)(P_i)$; there is no missing different factor in these
coordinates. The different is already encoded by the trace-dual
lattice.

\section{The lattice index and the sign of the degree inequality}
Put $O=\OK^d$. Since the fibre is generically reduced and split,
$C$ is a full order contained in $O$. At a finite prime, choose
an $\OK_{\mathfrak p}$-basis of $O$ and a matrix $M$ whose columns
form a basis of $C$. The trace form on $O$ is the identity, so
\[
 v_{\mathfrak p}(\operatorname{disc}C)
 =v_{\mathfrak p}(\det(M^{\mathsf T}M))
 =2\operatorname{length}(O/C).
\]
The trace dual has basis $M^{-\mathsf T}$, whence
\[
 O\subset C^\vee,\qquad
 \operatorname{length}(C^\vee/O)=\operatorname{length}(O/C).
\]
For any coordinate projection onto $g$ positions,
$C^\vee/O$ surjects onto $F/\OK^g$, where $F$ is the projected
lattice. With the ordinary Euclidean metrics on these coordinate
spaces,
\[
 \widehat{\deg}F=\log\#(F/\OK^g)
 \le\tfrac12\log\mathrm N\operatorname{disc}C.
\]
Here the degree of the trivial lattice with its standard metric is
zero. No number-field discriminant term occurs in this convention.

If $\phi:L\to M$ is an integral nonzero map of metrized lines,
the product formula gives
\[
 \widehat{\deg}L\le\widehat{\deg}M+
 \sum_{\sigma:K\hookrightarrow\mathbb C}\log\|\phi\|_\sigma.
\]
This verifies the direction of the degree inequality used in the
candidate argument.

\section{The archimedean estimate}
Above the disc $|t-2|<1/2$, the $d$ inverse branches are disjoint
and holomorphic. For $\omega=f_i(t)\,dt$ on the $i$th branch,
the mean-value inequality gives
\[
 |f_i(2)|^2\le\frac4\pi\int_{|t-2|<1/2}|f_i(t)|^2\,dx\,dy.
\]
The selected branch domains are disjoint in $X(\mathbb C)$.
Summing the inequalities over any selected coordinate set $I$ gives
the stronger operator estimate
\[
 \sum_{i\in I}|f_i(2)|^2\le\frac4\pi\|\omega\|_{L^2}^2.
\]
Thus every singular value of the projected evaluation map is at
most $2/\sqrt\pi$. Its determinant norm is at most
$(2/\sqrt\pi)^g$, as asserted. The choice of the coordinate set
over $K$ remains permissible at every embedding; the sheets may be
permuted, which does not change this estimate.

Riemann--Hurwitz gives $2g-2\le d-3$, so a differential vanishing
on the full unramified fibre is zero. Some $g$-coordinate projection
is therefore injective. Combining the preceding facts yields
\[
 \hF(\operatorname{Jac}X)\le
 \frac{\log\mathrm N\operatorname{disc}C}{2[K:\Q]}
 +g\log\frac2{\sqrt\pi}.
\]
Javanpeykar, Sections~2.1 and~2.3, explicitly identify the stated
$L^2$ determinant metric with the stable Jacobian Faltings height.

\section{Finite discriminant and base extension}
The trace determinant is a section of an invertible sheaf on the
regular base surface. Its divisor has no additional codimension-two
part. Hence its pullback to $t=2$ is computed by the intersections
with its horizontal and vertical divisors.

The horizontal coefficients are at most $d$ by characteristic-zero
tame ramification. Only the section $t=0$ meets $t=2$, with total
normalized intersection $\log2$. At a vertical divisor over
$\mathfrak p\mid p$, Proposition~4.1.3 of Javanpeykar states
\[
 r_j\le e_j-1+e_je_{\mathfrak p}
             \left\lfloor\frac{\log d}{\log p}\right\rfloor.
\]
The proposition explicitly permits imperfect residue fields; thus
it applies to the Gauss valuations with residue field
$k(\mathfrak p)(t)$. Summing with weights $f_j$ and using
$\sum_j e_jf_j=d$ gives exactly the candidate bound.

Finite etaleness outside the specified bad primes and branch
sections persists after base extension and normalization: the
base-changed etale algebra over a normal base is already normal.
Ramification of the enlarged number field causes no omitted term,
because $\sum_{\mathfrak p\mid p}e_{\mathfrak p}f_{\mathfrak p}
=[K:\Q]$ is used before normalizing the vertical contribution.

\section{The modular and isogeny steps}
For the displayed curve $E_t$, direct substitution gives
$\Delta=1728t^2/(t-1)^3$ and $j=1728t$. Exact-order-$N$ torsion
is finite etale away from $6N$ and the three branch sections.
The action of $H$ is free, and its quotient remains finite etale,
even at primes dividing $|H|$: this is a free action on a finite
etale scheme. Since $-1\in H$, the ambiguity from the generic
automorphism group $\{\pm1\}$ has been removed. This is the
claimed integral modular cover.

For the final height comparison let $L$ be the field, of the
bounded degree supplied by the inherited realization, over which
an actual surface $B\subset J=\operatorname{Jac}X_H$ is isogenous
to the desired surface $A$. R\'emond's Corollary~1.2, also valid
for stable heights, gives
\[
 \hF(B)\le\hF(J)+\kappa(\dim J-2).
\]
Gaudron--R\'emond's Theorem~1.9(1), applied with reference variety
$B$ of dimension two over $L$, bounds an isogeny $B\to A$ by a
fixed power of
\[
 [L:\Q]\max(1,\log[L:\Q],\hF(B)).
\]
The Faltings height change is at most half the logarithm of its
degree. This adds only
$O(\log[L:\Q]+\log\max(1,\hF(B)))$ to $\hF(B)$.
One must use this bounded field $L$, rather than unnecessarily
applying the isogeny theorem over the possibly much larger field
that splits the fibre and makes the curve semistable. Stable
heights allow the two extensions to be used independently.

\section*{Primary references checked}
\begin{itemize}
\item A.~Javanpeykar, with an appendix by P.~Bruin,
\emph{Polynomial bounds for Arakelov invariants of Belyi curves},
Algebra Number Theory 8 (2014), 89--140.
Sections~2.1, 2.3, 4.2 and Proposition~4.1.3.
\url{https://arxiv.org/pdf/1403.6404}.
\item G.~R\'emond, \emph{Propri\'et\'es de la hauteur de Faltings},
Ann. Sc. Norm. Super. Pisa Cl. Sci. 23 (2022), 1589--1596.
Corollary~1.2 and the explicit stable-height extension on page~2.
\url{https://www-fourier.univ-grenoble-alpes.fr/~remond/2476.pdf}.
\item E.~Gaudron and G.~R\'emond,
\emph{Nouveaux th\'eor\`emes d'isog\'enie},
M\'em. Soc. Math. France 176 (2023), Theorem~1.9(1).
\url{https://lmbp.uca.fr/~gaudron/art20.pdf}.
\end{itemize}
\end{document}
